\section{A product-preserving interval and the general module-category counit}

\subsection{F-SM40: the interval functor, rather than the resolution, can supply finiteness}

Let \(k\) be a commutative ring, let \(R\) be a small
\(k\)-linear category with projective scalar morphism modules,
and put \(\mathcal M=\operatorname{Mod}\!-\!R\).
Retain the scalar adjunctions, exact scalar functor, canonical
normalization, and exact resolution adjunction of BR-13403--05:
\[
 \begin{gathered}
 \varphi\dashv\kappa\dashv\psi,\qquad
 \kappa:\operatorname{Mod}(k)\longrightarrow\mathcal M
       \text{ exact},\qquad
 L\varphi(\kappa k)\longrightarrow k[0]\text{ a quasi-isomorphism},\\
 \varphi=D\otimes_R-,\qquad B_\bullet\longrightarrow D
       \text{ a projective left-}R\text{ resolution},\\
 \Lambda X=B_\bullet\otimes_R X,\qquad
 (GC)(x)=\operatorname{Hom}_{\operatorname{Ch}_{\ge0}(k)}
                                  (B_\bullet(x),C).
 \end{gathered}                                                \tag{SM40.1}
\]
The normalization map is the canonical derived counit, not
an unspecified isomorphism of objects.

Suppose there is a \(k\)-linear endofunctor \(K\) of
\(\mathcal M\) and a natural transformation
\(\alpha:\mathrm{Id}\to K\) such that
\[
 \alpha_X\text{ is monic and }L\varphi(KX)=0
                 \quad\text{for every }X\in\mathcal M.       \tag{SM40.2}
\]
Assume in addition that \(K\) preserves coproducts,
cokernels, and products.  In the original interval setup,
\(K=J\otimes_{\mathcal M}-\); thus this is an explicit
condition on the actual interval functor.  Exactness of
\(K\) is not required.

Then the ordinary counit \(\Lambda G(C)\to C\) is a
quasi-isomorphism for every connective scalar complex.
Consequently this very adjunction induces
\[
 \mathcal M[W^{-1}]\simeq D_{\ge0}(k),\qquad
                  W=\Lambda^{-1}(\mathrm{qis}).              \tag{SM40.3}
\]
More strongly, for arbitrary objects \(x_i\in R\) and
arbitrary scalar modules \(U_i\),
\[
 L\varphi\!\left(\prod_i C_{x_i}(U_i)\right)=0,
 \qquad C_x(U)(y)=\operatorname{Hom}_k(R(x,y),U).              \tag{SM40.4}
\]
No scalar dimension or cardinality bound, no finite generation
of \(R(x,y)\), and no \(\mathrm{FP}_\infty\) hypothesis on
\(D\) is imposed.

\begin{proof}
\emph{1. A finite-presentation criterion for tensor functors.}
We first prove the elementary ringoid version needed below.
For a left \(R\)-module \(L\), the functor
\(L\otimes_R-\) preserves products if and only if \(L\)
is finitely presented.  Here a finite presentation uses
finite coproducts of left representables \(R(a,-)\).

For the forward implication, index a family by all pairs
\(i=(a_i,\ell_i)\) with \(\ell_i\in L(a_i)\), and put
\(P_i=R(-,a_i)\).  Product preservation makes
\[
 L\otimes_R\prod_iP_i\longrightarrow
                \prod_i(L\otimes_RP_i)=\prod_iL(a_i)          \tag{SM40.5}
\]
surjective.  A preimage of the tuple \((\ell_i)_i\) is a
finite sum \(\sum_{j=1}^m\ell'_j\otimes w_j\), where
\(\ell'_j\in L(b_j)\) and
\(w_j\in\prod_iR(b_j,a_i)\).  Writing its coordinates
as \(r_{ji}:b_j\to a_i\) gives
\[
                  \ell_i=\sum_{j=1}^m L(r_{ji})(\ell'_j).
\]
Thus finitely many \(\ell'_j\) generate \(L\).  Choose a
finite coproduct \(F\) of left representables and an
epimorphism \(F\to L\), and write \(H\) for its kernel.

Now index a second family by every pair
\(i=(a_i,h_i)\) with \(h_i\in H(a_i)\), again using
\(P_i=R(-,a_i)\).  The tuple \((h_i)_i\), viewed in
\(\prod_i F(a_i)\), comes from a unique element
\(z\in F\otimes_R\prod_iP_i\): tensoring with a finite
coproduct of left representables is a finite coproduct of
evaluation functors, and therefore preserves products.
The image of \(z\) in \(L\otimes_R\prod_iP_i\) has zero
image under (SM40.5).  Injectivity of that comparison map
makes this image zero.  Right exactness of tensoring in
the left-module variable now writes
\[
                  z=\sum_{j=1}^s h'_j\otimes v_j
       \quad\text{in }F\otimes_R\prod_iP_i,
\]
with finitely many \(h'_j\in H(b_j)\).  Taking coordinates
shows that these \(h'_j\) generate every element of \(H\).
Consequently \(H\) is finitely generated and \(L\) is
finitely presented.

Conversely, tensor a finite presentation of \(L\) with
an arbitrary product of right modules.  The two finite
representable sums commute with the product.  Products
of scalar modules are exact, so taking the product of
the corresponding right exact presentations and comparing
cokernels proves the converse.  All index sets used above
are sets because \(R\) is small and its modules have
set-valued components.

\emph{2. Each evaluation of the interval has a finitely presented coefficient.}
For \(x\in R\) define a left \(R\)-module
\[
                   L_x(a)=(K(R(-,a)))(x).                   \tag{SM40.6}
\]
Its left action is induced by maps between right
representables.  There is a natural evaluation pairing
\[
                  L_x\otimes_R X\longrightarrow(KX)(x).
                                                               \tag{SM40.7}
\]
For an element of \(X(a)\), use its associated map
\(R(-,a)\to X\), apply \(K\), and evaluate at \(x\).
The \(k\)-linearity and naturality of this operation make
the pairing balanced.  It is an isomorphism on right
representables.  Since \(K\) preserves coproducts and
cokernels, presentations by coproducts of representables
make (SM40.7) an isomorphism for every \(X\).
This gives a direct proof of the required tensor
description, without an additional representation theorem.

Because \(K\) preserves products, (SM40.7) shows that
\(L_x\otimes_R-\) preserves products.  Step 1 makes
\(L_x\) finitely presented.  Under (SM40.7), the evaluated
natural transformation \(\alpha_X(x)\) is tensoring with
the left-module map
\[
                       u_x:R(x,-)\longrightarrow L_x        \tag{SM40.8}
\]
obtained by evaluating \(\alpha\) on right representables.
The assertion follows first on these representables and
then on their coproduct presentations.

\emph{3. The evaluated interval inclusions split.}
By (SM40.2), tensoring (SM40.8) with every right module
gives a monomorphism.  Thus \(u_x\) is a pure monomorphism
of left \(R\)-modules.  Its quotient \(V_x\) is finitely
presented.  Indeed, take a finite presentation of \(L_x\).
The map from the projective left representable \(R(x,-)\)
lifts to the finite free term of that presentation; adding
this one representable to its relation term gives a finite
presentation of \(V_x\).

A pure short exact sequence with finitely presented
quotient splits.  In this case the finite-equation argument
can be stated explicitly.  Choose finitely many generators
and relations for \(V_x\), and lift the generators to
\(L_x\).  Each relation on the lifts gives an element of
\(R(x,-)\).  Those finitely many linear equations, with
constants in \(R(x,-)\), have a solution in \(L_x\).
Purity gives a solution in \(R(x,-)\).  Subtracting that
solution from the chosen lifts makes all relations vanish
and defines a section \(V_x\to L_x\).  The finite-equation
criterion for purity is the one proved in F-SM24 by
tensoring the transposed finite presentation matrix.
We obtain a left-module retraction
\(r_x:L_x\to R(x,-)\) of \(u_x\).

Tensoring \(r_x\) with \(X\) now splits
\[
                     \alpha_X(x):X(x)\longrightarrow(KX)(x)
                                                               \tag{SM40.9}
\]
as a scalar map, naturally in \(X\).  The choices \(r_x\)
need not be natural in \(x\); no retraction \(K\to\mathrm{Id}\)
of functors into \(\mathcal M\) has been asserted.
This distinction is essential.

\emph{4. Arbitrary products of coinduced coefficients are killed.}
Put \(I=\prod_i C_{x_i}(U_i)\).  Coinduction gives
\[
 \operatorname{Hom}_{\mathcal M}(X,I)
             =\prod_i\operatorname{Hom}_k(X(x_i),U_i).       \tag{SM40.10}
\]
For every monomorphism whose evaluations split over \(k\),
the restriction map in (SM40.10) is surjective: extend each
coordinate by composing with its chosen scalar retraction.
Apply this to \(\alpha_I:I\to KI\), whose evaluations
split by Step 3.  The identity of \(I\) extends to a map
\(KI\to I\).  Hence \(I\) is a retract of \(KI\).
Equation (SM40.2) and functoriality of \(L\varphi\) show
that \(L\varphi(I)\) is a retract of zero, proving
(SM40.4).  Neither injectivity nor pure injectivity of
the coefficients \(U_i\) is involved.

\emph{5. The unreplaced counit.}
If \(P\) is an acyclic projective object of
\(\operatorname{Ch}_{\ge0}(k)\), BR-13402 writes it as
\[
           P=\prod_{n\ge1}D_n(Q_n),\qquad Q_n\text{ projective over }k.
                                                               \tag{SM40.11}
\]
The product is also the coproduct degree by degree, but
we use that \(G\), as a right adjoint, preserves the
product.  The disk adjunction gives
\[
            GP=\prod_{n\ge1}\operatorname{Hom}_k(B_{n-1}(-),Q_n).
                                                               \tag{SM40.12}
\]
Each projective left module \(B_{n-1}\) is a retract of
a coproduct of left representables.  Its scalar Hom
module is therefore a retract of a product of coinduced
modules \(C_x(Q_n)\).  Taking products of these retractions
makes (SM40.12) a retract of one product of such modules.
Step 4 and preservation of retracts give
\(L\varphi(GP)=0\).

For clarity, the canonical normalization in (SM40.1)
also handles every degree-zero projective scalar module.
The adjunction identifies \(G(Q[0])\) with \(\kappa Q\)
and its counit with the scalar derived counit.  The
functors \(\kappa\) and \(\Lambda\) preserve coproducts,
and scalar direct sums preserve quasi-isomorphisms.
The normalization for \(k\) therefore passes to free
modules and then to their projective retracts.

Every projective chain object is a finite biproduct of
such a degree-zero object and an acyclic projective
complex.  Its counit is consequently a quasi-isomorphism.
The exact-adjunction criterion F-SM17 extends this to
every target: resolve the target by projective chain
objects, apply the two exact functors \(\Lambda G\) and
\(\mathrm{Id}\), and compare their first-quadrant total
complexes.  The diagonals are finite, the augmentations
are quasi-isomorphisms, and the counit on every resolution
term is a quasi-isomorphism.  Naturality and the triangle
identities then descend the actual unit and counit to
inverse isomorphisms, proving (SM40.3).
\end{proof}

\paragraph{Exact remaining scope.}
F-SM40 imposes finiteness on the evaluated interval
coefficients \(L_x\), not on \(D\), its resolution terms,
or the scalar modules \(R(x,y)\).  Equivalently, among
coproduct- and cokernel-preserving interval functors,
Step 1 identifies the extra product-preservation condition
with finite presentation of all the \(L_x\).
The split-evaluation conclusion is stronger than the
objectwise purity used in F-SM24.

The handwritten bilinear tensor convention has not been
shown to imply that \(J\otimes-\) preserves products.
Thus this result removes another explicit class from
the general raw-counit gate; it does not declare that
gate closed under the bare original assumptions.
In particular, no passage from an arbitrary left adjoint
to product preservation is used.
